% ============================================================================
%  PLANTILLA IEEE — Artículo de Revista (Journal Paper)
%  Basada en IEEEtran v1.8b+ | Cumple con ieee-latex SKILL.md
%  Diferencias con conference: sin \IEEEoverridecommandlockouts, incluye
%  biografías de autores, usa \thanks para fechas y afiliaciones detalladas.
% ============================================================================

\documentclass[journal]{IEEEtran}

% ====================== PAQUETES ESENCIALES ==================================
\usepackage[utf8]{inputenc}
\usepackage[T1]{fontenc}
\usepackage{cite}
\usepackage{amsmath,amssymb,amsfonts}
\usepackage{graphicx}
\usepackage[caption=false]{subfig}
\usepackage{array}
\usepackage{booktabs}
\usepackage{multirow}
\usepackage{listings}
\lstset{
  basicstyle=\footnotesize\ttfamily,
  breaklines=true,
  frame=single,
  framesep=6pt,
  columns=fullflexible,
  keepspaces=true,
  showstringspaces=false,
  numbers=none,
  xleftmargin=8pt,
  xrightmargin=8pt,
  framexleftmargin=6pt,
  framexrightmargin=6pt,
  aboveskip=10pt,
  belowskip=6pt,
  captionpos=t,
  abovecaptionskip=10pt,
  belowcaptionskip=15pt,
  language=bash,
}
\usepackage{algorithm}
\usepackage{algorithmic}
\usepackage{url}
\usepackage[hidelinks]{hyperref}

\graphicspath{{figures/}}

% ====================== INICIO DEL DOCUMENTO =================================
\begin{document}

\title{Título del Artículo de Revista}

\author{Primer~A.~Autor,~\IEEEmembership{Member,~IEEE,}
  Segundo~B.~Autor,~and~Tercer~C.~Autor,~Jr.,~\IEEEmembership{Fellow,~IEEE}
  \thanks{Manuscript received Month Day, Year; revised Month Day, Year.}
  \thanks{P. A. Autor is with the Department of Engineering, University,
    City, Country (e-mail: primer@universidad.edu).}
  \thanks{S. B. Autor is with Company Name, City, Country.}
  \thanks{T. C. Autor is with Another University, City, Country.}
}

% Encabezado de página (opcional, nombre corto del artículo)
\markboth{Journal of \LaTeX\ Class Files, Vol.~XX, No.~X, Month~Year}%
{Autor \MakeLowercase{\textit{et al.}}: Título corto del artículo}

\maketitle

\begin{abstract}
  Abstract de 100 a 200 palabras para artículos regulares, máximo 50 palabras
  para correspondence/technote. Sin citas, ecuaciones ni caracteres especiales.
\end{abstract}

\begin{IEEEkeywords}
  Keyword uno, keyword dos, keyword tres, keyword cuatro, keyword cinco.
\end{IEEEkeywords}

% ============================================================================
\section{Introduction}
\label{sec:intro}

\IEEEPARstart{E}{l} primer párrafo de un artículo de journal IEEE usa una
letra capitular (drop cap) con el comando \verb|\IEEEPARstart{L}{etra}|.
El resto del artículo sigue la estructura estándar.

Contextualizar el problema y citar trabajos previos \cite{haykin2009neural}.

\section{Related Work}
\label{sec:related}

Revisión de la literatura \cite{smith2020example}, \cite{jones2021analysis}.

\section{Proposed Method}
\label{sec:method}

Descripción del método propuesto.

\begin{equation}
  y = \mathbf{w}^T \mathbf{x} + b
  \label{eq:model}
\end{equation}

\section{Experimental Results}
\label{sec:results}

Resultados en Table~\ref{tab:results} y Fig.~\ref{fig:result}.

\begin{table}[!t]
  \caption{Resultados Experimentales}
  \label{tab:results}
  \centering
  \begin{tabular}{lcc}
    \toprule
    \textbf{Método} & \textbf{Precisión (\%)} & \textbf{Tiempo (s)} \\
    \midrule
    Baseline    & 90.1 & 3.2 \\
    Propuesto   & \textbf{95.7} & 2.8 \\
    \bottomrule
  \end{tabular}
\end{table}

\begin{figure}[!t]
  \centering
  \includegraphics[width=\columnwidth]{example-image}
  \caption{Resultado experimental principal.}
  \label{fig:result}
\end{figure}

\section{Discussion}
\label{sec:discussion}

Interpretación de resultados.

\section{Conclusion}
\label{sec:conclusion}

Conclusiones y trabajo futuro.

% ---------- APÉNDICES (opcional) ---------------------------------------------
\appendices
\section{Demostración del Teorema 1}
\label{app:proof}
% Las ecuaciones en apéndices se numeran (A1), (A2), etc. automáticamente.

% ---------- AGRADECIMIENTOS --------------------------------------------------
\section*{Acknowledgment}
Los autores agradecen a\ldots

% ---------- BIBLIOGRAFÍA -----------------------------------------------------
\bibliographystyle{IEEEtran}
\bibliography{references}

% ---------- BIOGRAFÍAS (solo journal) ----------------------------------------
\begin{IEEEbiographynophoto}{Primer A. Autor}
  recibió el título de Ingeniero en Sistemas de la Universidad Ejemplo en 2020.
  Actualmente cursa la Maestría en Ciencias de la Computación en la misma
  universidad. Sus intereses de investigación incluyen almacenamiento
  distribuido y tolerancia a fallos.
\end{IEEEbiographynophoto}

\begin{IEEEbiographynophoto}{Segundo B. Autor}
  es Profesor Asociado en el Departamento de Ingeniería de la Otra Universidad.
  Sus líneas de investigación abarcan infraestructura computacional y sistemas
  operativos.
\end{IEEEbiographynophoto}

% Con foto (descomentar y ajustar):
% \begin{IEEEbiography}[{\includegraphics[width=1in,height=1.25in,clip,
%     keepaspectratio]{photos/tercer_autor.jpg}}]{Tercer C. Autor, Jr.}
%   Biografía con foto (1" × 1.25", 300 DPI, escala de grises).
% \end{IEEEbiography}

\end{document}
